\section{Worship, Reform, and the Name \emph{Gregorian}}

Few names are woven so tightly into the fabric of Western liturgy as Gregory's.
For over a millennium, Latin Christendom spoke of the ``Gregorian''
Sacramentary and ``Gregorian'' chant as though Pope Gregory the Great had
personally compiled the Roman Missal and composed its melodies. Separating
Gregory's verified liturgical acts from the later Carolingian and medieval
compilations that honored his name does not shrink his stature; it reveals how
the Roman liturgy actually grew.

\subsection{Gregory's verified liturgical work}

Direct sixth-century testimony and Gregory's own writings establish four
concrete areas of liturgical activity:

\begin{enumerate}
\item \textbf{Stational worship and preaching:} The headings and texts of the
\work{Homiliae XL in Evangelia} document Gregory presiding over Rome's
stational calendar across the city's titular churches and cemeteries, including
the three Christmas Masses celebrated on a single day (\work{Homiliae in
Evangelia} 8.1, delivered at Saint Mary Major) and Ember-day stations
(Hom.\ 20). Gregory of Tours (\work{Historia Francorum} X.1) preserves his
organization of the sevenfold penitential procession (\latin{litania
septiformis}) of 590.
\item \textbf{The Roman \emph{Kyrie} and the \emph{Pater noster}:}
In his letter to Bishop John of Syracuse (\work{Ep.}\ IX.12, quoted in
Fortescue 1922, pp.\ 234--235), Gregory distinguished Roman usage of the
\latin{Kyrie eleison} from Greek practice: ``We neither say nor have said
\latin{Kyrie eleison} as it is said by the Greeks \dots\ But in the daily
Masses we leave out some things which are generally said; we only say
\latin{Kyrie eleison} and \latin{Christe eleison}, that we should dwell rather
longer on these words of prayer.'' As Fortescue demonstrates (pp.\ 234--235),
Gregory did \emph{not} introduce the \latin{Kyrie} into the Roman Mass---the
Council of Vaison (529) already proves it older---but regulated its distinctive
Roman alternation with \latin{Christe eleison} and simplified daily usage.
Alongside that regulation, consistent Roman liturgical testimony credits
Gregory with inserting the Lord's Prayer (\latin{Pater noster}) before
Communion after the Canon (Fortescue 1922, p.\ 172).
\item \textbf{The \emph{Hanc igitur} of the Roman Canon:} Consistent Roman
tradition (recorded by Bede, the \work{Liber Pontificalis}, and Fortescue
1922, p.\ 172) credits Gregory with adding the haunting petition for earthly
peace and eternal salvation to the \latin{Hanc igitur} inside the Roman Canon:
\latin{diesque nostros in tua pace disponas, atque ab aeterna damnatione nos
eripi, et in electorum tuorum iubeas grege numerari}---``and dispose our days
in thy peace, and bid us to be delivered from eternal damnation and numbered in
the flock of thine elect'' (\work{Hadrianum} Canon, ed.\ Wilson 1915).
Written in a Rome surrounded by Lombard armies, that clause remained the last
substantive textual addition to the core of the Roman Canon for more than a
millennium (Fortescue 1922, p.\ 172, citing Benedict XIV, \work{De ss.\ Missae
Sacrificio} II.xii.12).
\item \textbf{Unity in faith, freedom in local rite:} Gregory never pursued a
rigid liturgical uniformity across Christendom. Advising Leander of Seville on
whether Spanish Catholics should use single or triple immersion in baptism so
as to distinguish themselves from Arians, he stated the guiding principle
plainly: ``where there is one faith, a diversity of usage does no harm to holy
Church'' (\work{Ep.}\ I.43).
\end{enumerate}

\subsection{The ``Gregorian'' Sacramentary (\emph{Hadrianum})}

Nearly two centuries after Gregory's death, Charlemagne asked Pope Hadrian I
(772--795) for an authentic Roman Mass book to guide liturgical reform across
the Frankish realm. Between 784 and 791, Hadrian sent north a papal
sacramentary drawn from the Lateran archive (\latin{ex authentico libro
bibliothecae cubiculi}), bearing the title \latin{Liber Sacramentorum de
circulo anni expositus a sancto Gregorio papa Romano editus}---``The Book of
Sacraments for the cycle of the year, set forth by Saint Gregory, Roman Pope''
(Hadrian I's letter and the manuscript rubric printed in Migne PL 78, cols.\
23--26, and edited critically by H.\ A.\ Wilson, 1915, pp.\ xxi--xxii).

The text of the \work{Hadrianum} (Wilson 1915) shows at once why the whole book
cannot be Gregory's own sixth-century manuscript. It
already incorporates seventh- and eighth-century Roman developments, including
the feasts of the Nativity and Assumption of the Blessed Virgin, Lenten
Thursday Masses, and---most tellingly---a complete Mass for Saint Gregory's
own feast day on 12 March: \latin{IIII Idus Marcias, id est XII die mensis
Marcii, Natale Sancti Gregorii Papae} (Wilson 1915, p.\ 25).

Even more revealing is what happened once the book reached Francia. Because
Hadrian had sent a specialized papal stational book designed for the pope's
own solemn celebrations in Rome, it lacked ordinary Sunday formularies,
votive Masses, blessings, and parochial rites that Frankish priests needed
every week. A Carolingian compiler---attributed by the eleventh-century
\work{Micrologus} of Bernold of Constance and by Wilson (1915, pp.\ xxi--xxii)
to Alcuin, though later twentieth-century scholarship has largely favored
Saint Benedict of Aniane---assembled a large Supplement from older Roman and
Gallican sources and placed between the two halves the famous \latin{Hucusque}
preface (Wilson 1915, pp.\ 145--146). With remarkable scholarly honesty, the
\latin{Hucusque} writer drew a boundary line right down the middle of the
codex: ``Up to this point (\latin{Hucusque}),'' he explained, stands the book
of Blessed Pope Gregory, except for passages marked in the margin with
obelisks (\latin{virgulis antepositis})---explicitly singling out the Mass of
Saint Gregory himself as ``inserted without doubt into his work by his
successors out of love and veneration for him'' (Wilson 1915, p.\ 145). The
\work{Hadrianum} thus preserves a genuine Roman papal nucleus rooted in
Gregory's age and reworked across the seventh and eighth centuries, fused in
Charlemagne's empire with a Frankish Supplement. It honors Gregory as the
great patron of the Roman rite without being a single volume penned by his
hand.

\subsection{Gregorian chant and the \emph{Antiphonarius}}

A parallel history governs the music that bears Gregory's name. Nearly three
hundred years after Gregory's death, John the Deacon (\work{Vita S. Gregorii}
II.6, c.\ 873--875, quoted in the Maurist \work{Ad Antiphonarium S. Gregorii
Praefatio}, Migne PL 78, col.\ 649) wrote that Gregory, ``after the manner of
most wise Solomon, for the sweet compunction of music, most usefully compiled
an Antiphonary patchwork (\latin{Antiphonarium centonem}),'' established the
\latin{schola cantorum} (singing school) at Rome, and taught choristers from
his sickbed with a rod (\latin{flagello}) still shown to ninth-century
visitors. A luxurious ninth-century Compi\`egne manuscript carried the famous
prologue \latin{Gregorius praesul} crediting him with renewing the monuments of
earlier Fathers and composing a musical handbook for the \latin{schola
cantorum} (Migne PL 78, col.\ 650).

Yet even the seventeenth-century Maurist editors and Pierre de Goussainville
(Migne PL 78, cols.\ 649--651) sounded a clear note of caution. The word
\latin{cento} means a compilation stitched from earlier materials, not an
original composition; ecclesiastical chant was far older than Gregory; and
John the Deacon wrote three centuries after the fact. Moreover, John himself
(\work{Vita} II.9, quoted in Migne PL 78, col.\ 650) unwittingly points to the
real historical crucible of what became ``Gregorian chant'': the fierce
eighth-century collision between Roman and Gallican singing under Pepin and
Charlemagne, when Frankish cantors studied Roman practice and carried it back
to Metz and across Gaul. Modern musicology confirms that the standardized
repertoire of Latin plainchant crystallized in that eighth- and ninth-century
Frankish-Roman synthesis---drawing on Old Roman roots, reshaped in Carolingian
cathedrals and abbeys, and placed under the prestigious banner of Pope Saint
Gregory, whose care for Roman worship and whose English mission had first
carried Roman psalmody beyond the Alps.
